\documentclass[11pt]{letter}
\usepackage[letterpaper, left=0.8in, right=0.8in, top=0.65in, bottom=0.65in]{geometry}
\usepackage[hidelinks]{hyperref}
\usepackage{parskip}
\pagestyle{empty}

\begin{document}

\begin{center}
{\Large \textbf{Aryan Miriyala}}\\
\href{mailto:aryanmiriyala@gmail.com}{aryanmiriyala@gmail.com} $|$ +1 419-315-0444 $|$
\href{https://www.linkedin.com/in/aryan-miriyala-7788a21b6/}{linkedin.com/in/aryan-miriyala} $|$
\href{https://github.com/aryanmiriyala}{github.com/aryanmiriyala}
\end{center}

\vspace{0.15in}

Dear MeritFirst Team,

What stands out to me about MeritFirst is the focus on seeing what early-career engineers can actually build. That resonates with me because most of my strongest experience has come from taking messy, practical problems and turning them into working systems: secure internal tools at SmartSolve, production data automation at AAIS, reliability fixes for a healthcare workflow platform at APKD, and personal projects where I owned the full path from idea to implementation. I like environments where the work sample matters because it makes the conversation more honest: can I understand the problem, write working code, debug it, explain tradeoffs, and keep improving?

I would bring a broad builder's foundation across Python, TypeScript, C++, SQL, React, Next.js, Node.js, PostgreSQL, MongoDB, Docker, and AWS. I have built full-stack products like RocketGrader and Travel Health Advisor, automated multi-terabyte production data workflows with Python/PySpark/AWS Glue, implemented RBAC/JWT and auth patterns, and pushed through a self-adaptive C++/OpenMP scheduling project that taught me how to reason from measurement instead of assumptions. I am interested in MeritFirst because it points toward the kind of engineering career I want: shipping real code, learning quickly from strong engineers, taking feedback seriously, and proving fit through the quality of what I build.

\end{document}
